Optical and electrophysiological neural recordings have exponentially increased in scale over the past several decades~\parencite{stevenson2011advances,urai2022large}.
In particular, population-level neural recordings can now be obtained across many subjects, brain regions, and repeated behavioral sessions, but summarizing these multi-session datasets is a formidable challenge.
Here, we used shape metrics to show that neural population geometry is neither uniform across the brain nor purely idiosyncratic across animals, but structured at both scales: organized by anatomy across regions and by behavior and genotype across individuals.

It is now widely appreciated that many task and behavioral variables can be decoded from nearly any recorded brain region \parencite{steinmetz2019distributed,stringer2019spontaneous,musall2019,angelaki2025brain, Lloreda2025,rosen2026distributed}.
This has been read as evidence for a highly distributed code(\textcite{hayden2026rethinking,fine2021whole,Lloreda2025}), but decodability of a single variable is a coarse summary of a population's activity.
Our results suggest a complementary view: when regions are compared at the level of population geometry rather than decodability, a clear anatomical organization re-emerges, with shape distance recovering a hierarchy across dozens of regions without any anatomical labels supplied to the algorithm (\textit{Fig.} \ref{fig:allen}, \ref{fig:miller}).

Further, individual differences in behavior were mirrored by individual differences in population geometry, extending to full cohorts a link that prior work had established only for pairs of subjects.
For instance, recent work by \textcite{fascianelli2024neural} performed a detailed analysis of two nonhuman primate subjects and found correlations between their behavioral differences and neural geometries.
Shape metrics provide a rigorous framework to scale up these investigations to larger experimental cohorts (e.g. across $K = 56$ regions in \textit{Fig.} \ref{fig:allen} or $K = 56$ mice in \textit{Fig.} \ref{fig:ibl}). Moreover, the same approach extends beyond spiking data. Applied to human EEG, shape distances reveal reliable subject-specific geometry, and show that self-supervised pretraining of a foundation model \parencite{cui2024,wang2025} amplifies rather than averages out these idiosyncrasies \parencite{marty2026}.

A further consequence of defining similarity through explicit alignment transformations is that the alignments themselves are useful objects.
Once a collection of neural systems has been aligned to a common reference (by generalized Procrustes analysis, in our case) one obtains a shared coordinate system in which all subjects can be jointly analyzed.
We exploited this to decode head direction in a held-out mouse from a decoder fit entirely on other mice (\textit{Fig.} \ref{fig:hdcells}e).
This links our framework directly to hyperalignment \parencite{haxby2020hyperalignment, Chen2021} and to latent-space alignment methods developed for across-session \parencite{degenhart2020stabilization, gallego2020long} and across-subject decoding \parencite{safaie2023preserved}. 
Distance and alignment are two aspects of the same computation: the distance quantifies how much structure two systems share, while the alignment supplies the map that makes shared structure usable.
Similarity measures that return only a scalar score do not furnish this map (\hl{who to cite}).

Shape metrics are not the only family of distance functions that satisfy symmetry and the triangle inequality.
Slight modifications to CKA scores and CCA-based similarity scores result in proper metric spaces~\parencite{Williams2021}.
Distance-based RSA scores can also be turned in proper metrics, as can be seen from recent work showing equivalence between RSA and CKA~\parencite{Williams2024}. These alternatives are complementary to ours. The contribution here is not the existence of a metric per se, but the demonstration of what a proper metric space buys on neural data: clustering, regression, and nearest-neighbor prediction across large collections of recordings.
 For simplicity, we utilized the Procrustes distance as a running example throughout most of our analyses.
Our practical experience suggests that the Procrustes distance is reasonable default choice, because it is resistant to overfitting and has no hyperparameters that need to be tuned, but practitioners should consider a broader set of distance functions depending on their analysis goals (\hl{who to cite}).

A distinct approach uses decoders themselves to characterize geometry. For example, cross-condition generalization (e.g. training a classifier in one condition and testing it in another) has been used to quantify how abstract or disentangled a representation is \parencite{bernardi,saez,barbosa2023,courellis2024, fascianelli2024neural,posani2026,nogueira2026}.
Decoding therefore probes geometry too, but through a different lens. Because a decoder targets a specific variable, it answers a question that shape distances cannot: whether a particular piece of information is present in a population, and in a format that a downstream reader could exploit.
A shape distance instead summarizes the geometry as a whole, which is the appropriate summary when the question concerns the structure of a representation rather than the accessibility of one variable within it.
The two are naturally used together and we indeed relied on decoding throughout this work to validate and interpret the structure that shape distances revealed.

The potential of the shape metrics framework to study neural geometry was not fully explored here. For instance, we have focused on analyzing trial-averaged neural responses, but the scale and orientation of trial-to-trial noise is important to theories of neural coding \parencite{moreno2014information}.
Furthermore, there is substantial interest in moving beyond geometry to understand neural circuits as dynamical systems \parencite{vyas2020computation}.
The framework we've outlined can be extended to address these issues \parencite{duong2023representational,ostrow2024beyond,nejatbakhsh2024comparing, cayco2026}.
However, these extensions remain a relatively active area of research and require hyper parameter fine-tuning.

Overall, our results reveal that neural population codes carries structure at multiple scales: an anatomical hierarchy across regions recovered without supervision, and a correspondence between population geometry and behavior.
As recording technologies continue to advance, generating increasingly large and complex neural datasets, shape metrics have the potential to become a fundamental tool for neuroscience, enabling researchers to systematically analyze large collections of neural systems at the level of population geometry.